\begin{abstract}
Large vision foundation models pretrained on massive datasets have demonstrated remarkable performance across various visual tasks. However, when applied to domain-specific applications such as semi-supervised medical image segmentation under domain shift, their strong prior knowledge may lead to overconfident predictions even when some of them are incorrect. These inaccurate pseudo-labels can accumulate during training, reducing the effectiveness of unlabeled data and limiting segmentation performance. In this work, we reproduce the SynFoC (Synergistic Training for Foundation and Conventional Models) framework, which jointly trains a foundation model and a conventional segmentation model to alleviate this issue. The conventional model, trained from scratch, helps correct the high-confidence errors made by the foundation model, while the foundation model provides reliable pseudo-labels that guide the conventional model during the early stages of training. The framework further employs consensus-divergence consistency regularization to encourage effective collaboration between the two models and improve training stability. We evaluate the reproduced implementation on four public multi-domain medical image segmentation datasets and report the experimental results obtained in our reproduction study. The implementation follows the original code released by the paper's authors at \url{https://github.com/MQinghe/SynFoC}.
\end{abstract}

\begin{IEEEkeywords}
medical image segmentation, semi-supervised learning, mixed-domain learning, foundation model, MedSAM, U-Net
\end{IEEEkeywords}
